An \code{EventBreakpoint} triggers when the target process encounters
a breakpoint
inserted by ProcControlAPI.

Similar to EventStop, EventBreakpoint will explicitly move the thread or process
to a stopped state. Returning cbDefault from a callback function that has
received EventBreakpoint will leave the target process in a stopped state
rather than restore it to the pre-event state.

Associated EventType Code: \code{Breakpoint}

\begin{apient}
  Dyninst::Address getAddress() const
\end{apient}

\apidesc{
  This function returns the address at which the breakpoint was hit.
}

\begin{apient}
  void getBreakpoints(std::vector<Breakpoint::const_ptr> &b) const
\end{apient}

\apidesc{
  This function returns a vector of pointers to the Breakpoints that were hit.
  Since it is possible to insert multiple Breakpoints at the same location, it
  is possible for this function to return more than one Breakpoint.
}
